% GENERATED FILE. Edit canonical lessons, then run npm run generate:books.
% canonical-source-hash: fnv1a64:70aaadb6160194b0
% canonical-lessons: GU-C408-saphar, GU-C408-phero, GU-C408-ato, GU-C408-visamo, GU-C408-tahelvu, GU-R408-review-401-408

\chapter{Trips and Rounds}
\label{ch:gu-trips-and-rounds}
\hlchaptermodality{408}

\begin{chapteropening}
\textbf{By the end of this chapter:} \emph{I can talk about trips, rounds and rests on the way.}

Complete the last lesson of chapter 408: I can talk about trips, rounds and rests on the way.
\end{chapteropening}

\section[\texorpdfstring{\gu{સફર} (saphar)}{સફર (saphar)}]{\gu{સફર} (saphar) — a voyage; a long journey}

\label{lesson:GU-C408-saphar}



\begin{quote}
Before the new one: say the Gujarati for on foot, then the Gujarati for to pine for home.
\end{quote}

\subsection*{You'll want to know: \gu{સફર}}
\textbf{\gu{સફર}} — \emph{saphar} — \textquotedblleft{}a voyage; a long journey\textquotedblright{}.

\begin{quote}
\gu{દરિયાની} \gu{સફર} \gu{મને} \gu{ઘણી} \gu{ગમી}.
\end{quote}

\begin{quote}
\emph{dariyānī saphar mane ghaṇī gamī.}
\end{quote}

\begin{quote}
I really enjoyed the sea voyage.
\end{quote}

A feminine noun.

\begin{cousinweb}
\textbf{\gu{સફર}} (\emph{saphar}) is usually traced to Arabic \emph{safar}, \textquotedblleft{}a journey\textquotedblright{}, which came through Persian; \gu{મુસાફર}, a traveller, is built on the same word.
\end{cousinweb}

\begin{grammarlens}[title={What you've built}]
One more word for getting around.
\end{grammarlens}

\subsection*{Your turn}
\begin{itemize}
\raggedright
  \item \emph{Say it:} \emph{saphar}
  \item \emph{Say it:} \emph{saphar}, once more
  \item \emph{Say it:} \emph{dariyānī saphar mane ghaṇī gamī.} — I really enjoyed the sea voyage.
  \item \emph{Recall:} say the Gujarati for on foot, then the Gujarati for to pine for home, then say \emph{saphar} again
\end{itemize}

\begin{tcolorbox}[breakable,colback=teal!4,colframe=teal!35!black,title={Before you move on}]
What does \gu{સફર} mean? (\textquotedblleft{}A voyage; a long journey\textquotedblright{}.) Say it once more.
\end{tcolorbox}
\section[\texorpdfstring{\gu{ફેરો} (phero)}{ફેરો (phero)}]{\gu{ફેરો} (phero) — a round trip; a round}

\label{lesson:GU-C408-phero}



\begin{quote}
Before the new one: say the Gujarati for to pine for home, then the Gujarati for a voyage.
\end{quote}

\subsection*{You'll want to know: \gu{ફેરો}}
\textbf{\gu{ફેરો}} — \emph{phero} — \textquotedblleft{}a round trip; a round\textquotedblright{}.

\begin{quote}
\gu{આ} \gu{બસ} \gu{દિવસમાં} \gu{ચાર} \gu{ફેરા} \gu{કરે} \gu{છે}.
\end{quote}

\begin{quote}
\emph{ā bas divasmā̃ cār pherā kare che.}
\end{quote}

\begin{quote}
This bus does four round trips a day.
\end{quote}

A masculine noun.

\begin{cousinweb}
\textbf{\gu{ફેરો}} (\emph{phero}) goes back to \emph{*phēra}, \textquotedblleft{}a turn, a turning\textquotedblright{}, a form CDIAL reconstructs beside \emph{*phirati}, \textquotedblleft{}moves, turns\textquotedblright{}; CDIAL lists the Gujarati word (CDIAL 9078). \gu{ફરવું}, to go around, is its relative.
\end{cousinweb}

\begin{grammarlens}[title={What you've built}]
One more word for getting around.
\end{grammarlens}

\subsection*{Your turn}
\begin{itemize}
\raggedright
  \item \emph{Say it:} \emph{phero}
  \item \emph{Say it:} \emph{phero}, once more
  \item \emph{Say it:} \emph{ā bas divasmā̃ cār pherā kare che.} — This bus does four round trips a day.
  \item \emph{Recall:} say the Gujarati for to pine for home, then the Gujarati for a voyage, then say \emph{phero} again
\end{itemize}

\begin{tcolorbox}[breakable,colback=teal!4,colframe=teal!35!black,title={Before you move on}]
What does \gu{ફેરો} mean? (\textquotedblleft{}A round trip; a round\textquotedblright{}.) Say it once more.
\end{tcolorbox}
\section[\texorpdfstring{\gu{આંટો} (ā̃ṭo)}{આંટો (ā̃ṭo)}]{\gu{આંટો} (ā̃ṭo) — a turn; a short walk}

\label{lesson:GU-C408-ato}



\begin{quote}
Before the new one: say the Gujarati for a voyage, then the Gujarati for a round trip.
\end{quote}

\subsection*{You'll want to know: \gu{આંટો}}
\textbf{\gu{આંટો}} — \emph{ā̃ṭo} — \textquotedblleft{}a turn; a short walk\textquotedblright{}.

\begin{quote}
\gu{સાંજે} \gu{અમે} \gu{બગીચામાં} \gu{આંટો} \gu{માર્યો}.
\end{quote}

\begin{quote}
\emph{sā̃je ame bagīcāmā̃ ā̃ṭo māryo.}
\end{quote}

\begin{quote}
In the evening we took a turn in the garden.
\end{quote}

A masculine noun.

\begin{cousinweb}
\textbf{\gu{આંટો}} (\emph{ā̃ṭo}) is a tadbhava of Sanskrit \emph{āvarta}, \textquotedblleft{}a whirlpool, a turning\textquotedblright{}; CDIAL lists the Gujarati word, \textquotedblleft{}a turn, a twist\textquotedblright{} (CDIAL 1418). \gu{આંટો} \gu{મારવો} is to go for a short walk.
\end{cousinweb}

\begin{grammarlens}[title={What you've built}]
One more word for getting around.
\end{grammarlens}

\subsection*{Your turn}
\begin{itemize}
\raggedright
  \item \emph{Say it:} \emph{ā̃ṭo}
  \item \emph{Say it:} \emph{ā̃ṭo}, once more
  \item \emph{Say it:} \emph{sā̃je ame bagīcāmā̃ ā̃ṭo māryo.} — In the evening we took a turn in the garden.
  \item \emph{Recall:} say the Gujarati for a voyage, then the Gujarati for a round trip, then say \emph{ā̃ṭo} again
\end{itemize}

\begin{tcolorbox}[breakable,colback=teal!4,colframe=teal!35!black,title={Before you move on}]
What does \gu{આંટો} mean? (\textquotedblleft{}A turn; a short walk\textquotedblright{}.) Say it once more.
\end{tcolorbox}
\section[\texorpdfstring{\gu{વિસામો} (visāmo)}{વિસામો (visāmo)}]{\gu{વિસામો} (visāmo) — a rest on the way; a resting place}

\label{lesson:GU-C408-visamo}



\begin{quote}
Before the new one: say the Gujarati for a round trip, then the Gujarati for a short walk.
\end{quote}

\subsection*{You'll want to know: \gu{વિસામો}}
\textbf{\gu{વિસામો}} — \emph{visāmo} — \textquotedblleft{}a rest on the way; a resting place\textquotedblright{}.

\begin{quote}
\gu{ચઢાણ} \gu{પછી} \gu{અમે} \gu{મંદિર} \gu{પાસે} \gu{વિસામો} \gu{લીધો}.
\end{quote}

\begin{quote}
\emph{caḍhāṇ pachī ame mandir pāse visāmo līdho.}
\end{quote}

\begin{quote}
After the climb we stopped for a rest by the temple.
\end{quote}

A masculine noun.

\begin{cousinweb}
\textbf{\gu{વિસામો}} (\emph{visāmo}) is a tadbhava of Sanskrit \emph{viśrāma}, \textquotedblleft{}rest\textquotedblright{}, and so \textquotedblleft{}a resting place\textquotedblright{}; CDIAL lists the Gujarati word (CDIAL 11956).
\end{cousinweb}

\begin{grammarlens}[title={What you've built}]
One more word for getting around.
\end{grammarlens}

\subsection*{Your turn}
\begin{itemize}
\raggedright
  \item \emph{Say it:} \emph{visāmo}
  \item \emph{Say it:} \emph{visāmo}, once more
  \item \emph{Say it:} \emph{caḍhāṇ pachī ame mandir pāse visāmo līdho.} — After the climb we stopped for a rest by the temple.
  \item \emph{Recall:} say the Gujarati for a round trip, then the Gujarati for a short walk, then say \emph{visāmo} again
\end{itemize}

\begin{tcolorbox}[breakable,colback=teal!4,colframe=teal!35!black,title={Before you move on}]
What does \gu{વિસામો} mean? (\textquotedblleft{}A rest on the way; a resting place\textquotedblright{}.) Say it once more.
\end{tcolorbox}
\section[\texorpdfstring{\gu{ટહેલવું} (ṭahelvũ)}{ટહેલવું (ṭahelvũ)}]{\gu{ટહેલવું} (ṭahelvũ) — to walk up and down; to stroll}

\label{lesson:GU-C408-tahelvu}



\begin{quote}
Before the new one: say the Gujarati for a short walk, then the Gujarati for a rest on the way.
\end{quote}

\subsection*{You'll want to know: \gu{ટહેલવું}}
\textbf{\gu{ટહેલવું}} — \emph{ṭahelvũ} — \textquotedblleft{}to walk up and down; to stroll\textquotedblright{}.

\begin{quote}
\gu{જમ્યા} \gu{પછી} \gu{દાદા} \gu{આંગણામાં} \gu{ટહેલે} \gu{છે}.
\end{quote}

\begin{quote}
\emph{jamyā pachī dādā ā̃gṇāmā̃ ṭahele che.}
\end{quote}

\begin{quote}
After eating, grandfather walks up and down the courtyard.
\end{quote}

Past (he …): \gu{તે} \gu{આંગણામાં} \gu{ટહેલ્યો}, \emph{te ā̃gṇāmā̃ ṭahelyo}.

\begin{cousinweb}
\textbf{\gu{ટહેલવું}} (\emph{ṭahelvũ}) goes back to \emph{*ṭahall}, \textquotedblleft{}to walk up and down\textquotedblright{}, a form CDIAL reconstructs; CDIAL lists the Gujarati verb, in the sense of going round from house to house, and Hindi \emph{ṭahalnā}, \textquotedblleft{}to take a walk\textquotedblright{} (CDIAL 5453).
\end{cousinweb}

\begin{grammarlens}[title={What you've built}]
One more word for getting around.
\end{grammarlens}

\subsection*{Your turn}
\begin{itemize}
\raggedright
  \item \emph{Say it:} \emph{ṭahelvũ}
  \item \emph{Say it:} \emph{ṭahelvũ}, once more
  \item \emph{Say it:} \emph{jamyā pachī dādā ā̃gṇāmā̃ ṭahele che.} — After eating, grandfather walks up and down the courtyard.
  \item \emph{Recall:} say the Gujarati for a short walk, then the Gujarati for a rest on the way, then say \emph{ṭahelvũ} again
\end{itemize}

\begin{tcolorbox}[breakable,colback=teal!4,colframe=teal!35!black,title={Before you move on}]
What does \gu{ટહેલવું} mean? (\textquotedblleft{}To walk up and down; to stroll\textquotedblright{}.) Say it once more.
\end{tcolorbox}
\section[(dialogue)]{Looking back: from A Flat Tyre to Trips and Rounds}

\label{lesson:GU-R408-review-401-408}



\begin{quote}
Say the Gujarati for a tyre and a puncture.
\end{quote}

\subsection*{Your turn}
Say these aloud:

\begin{itemize}
\raggedright
  \item a tyre, a puncture, a truck, to drag, to overturn
  \item a chhakda, a small goods van, a bullock cart, a bus conductor, to drive a cart
  \item a highway, a mile, a kilometre, to branch off, to crash (into)
  \item a pavement, a trail, narrow, an uphill climb, to jump over
  \item a flight, a pilot, an airport gate, a suitcase, a circle in the air
  \item a visa, a boundary, a customs duty, an official fee, to empty out
  \item abroad, a homeland, a dollar, on foot, to pine for home
  \item a voyage, a round trip, a short walk, a rest on the way, to walk up and down
\end{itemize}

\begin{tcolorbox}[breakable,colback=teal!4,colframe=teal!35!black,title={Before you move on}]
A tyre? (\textbf{\gu{ટાયર}}, \emph{ṭāyar}.) A puncture? (\textbf{\gu{પંક્ચર}}, \emph{pankcar}.) A truck? (\textbf{\gu{ટ્રક}}, \emph{ṭrak}.) To drag? (\textbf{\gu{ઘસડવું}}, \emph{ghasaḍvũ}.) To overturn? (\textbf{\gu{ઉથલવું}}, \emph{uthalvũ}.) A chhakda? (\textbf{\gu{છકડો}}, \emph{chakḍo}.) A small goods van? (\textbf{\gu{ટેમ્પો}}, \emph{ṭempo}.) A bullock cart? (\textbf{\gu{ગાડું}}, \emph{gāḍũ}.) A bus conductor? (\textbf{\gu{કંડક્ટર}}, \emph{kanḍakṭar}.) To drive a cart? (\textbf{\gu{હાંકવું}}, \emph{hā̃kvũ}.) A highway? (\textbf{\gu{ધોરીમાર્ગ}}, \emph{dhorīmārg}.) A mile? (\textbf{\gu{માઈલ}}, \emph{māīl}.) A kilometre? (\textbf{\gu{કિલોમીટર}}, \emph{kilomīṭar}.) To branch off? (\textbf{\gu{ફંટાવું}}, \emph{phanṭāvũ}.) To crash (into)? (\textbf{\gu{ટકરાવું}}, \emph{ṭakrāvũ}.) A pavement? (\textbf{\gu{ફૂટપાથ}}, \emph{phūṭpāth}.) A trail? (\textbf{\gu{કેડી}}, \emph{keḍī}.) Narrow? (\textbf{\gu{સાંકડું}}, \emph{sā̃kḍũ}.) An uphill climb? (\textbf{\gu{ચઢાણ}}, \emph{caḍhāṇ}.) To jump over? (\textbf{\gu{ટપવું}}, \emph{ṭapvũ}.) A flight? (\textbf{\gu{ઉડાન}}, \emph{uḍān}.) A pilot? (\textbf{\gu{પાયલટ}}, \emph{pāylaṭ}.) An airport gate? (\textbf{\gu{ગેટ}}, \emph{geṭ}.) A suitcase? (\textbf{\gu{સૂટકેસ}}, \emph{sūṭkes}.) A circle in the air? (\textbf{\gu{ચકરાવો}}, \emph{cakrāvo}.) A visa? (\textbf{\gu{વિઝા}}, \emph{vijhā}.) A boundary? (\textbf{\gu{સીમા}}, \emph{sīmā}.) A customs duty? (\textbf{\gu{જકાત}}, \emph{jakāt}.) An official fee? (\textbf{\gu{ફી}}, \emph{phī}.) To empty out? (\textbf{\gu{ઠાલવવું}}, \emph{ṭhālavvũ}.) Abroad? (\textbf{\gu{પરદેશ}}, \emph{pardeś}.) A homeland? (\textbf{\gu{વતન}}, \emph{vatan}.) A dollar? (\textbf{\gu{ડોલર}}, \emph{ḍolar}.) On foot? (\textbf{\gu{પગપાળા}}, \emph{pagpāḷā}.) To pine for home? (\textbf{\gu{હિજરાવું}}, \emph{hijrāvũ}.) A voyage? (\textbf{\gu{સફર}}, \emph{saphar}.) A round trip? (\textbf{\gu{ફેરો}}, \emph{phero}.) A short walk? (\textbf{\gu{આંટો}}, \emph{ā̃ṭo}.) A rest on the way? (\textbf{\gu{વિસામો}}, \emph{visāmo}.) To walk up and down? (\textbf{\gu{ટહેલવું}}, \emph{ṭahelvũ}.)
\end{tcolorbox}
